\section{Contractions in max norm}\label{sec:contractions}

At a point \(x\) far from a fixed point \(v\), a contraction in max
norm must move each coordinate of greatest error toward \(v\).
The signs of these movements give exactly the kind of partial
information used by Theorem~\ref{thm:prices}. We use them to locate
\(v\) with a running time polynomial in the number of bits of accuracy.

Let \(C=[-1,1]^n\), \(n\ge1\). A self-map \(f:C\to C\) with
\[
 \|f(x)-f(y)\|_\infty\le (1-\gamma)\|x-y\|_\infty,\qquad
 0<\gamma\le1,
\]
has a unique fixed point \(v\)~\cite{Ban}. Direct iteration may
require order \(1/\gamma\) evaluations to halve the error.
We assume an oracle that accepts a dyadic point \(x\in C\) and a
precision \(m\). It returns a dyadic point of \(C\) with \(m\) bits
after the binary point and sup norm error at most \(2^{-m}\) from
\(f(x)\). Dimensions and accuracy parameters are given in unary.

\begin{theorem}[Approximation in max norm]\label{thm:contraction}
For \(\gamma=2^{-b}\)
we can, given this access to the contraction \(f\)
and nonnegative integers \(b,k\), return
a dyadic \(x\in C\) with \(\|x-v\|_\infty\le 2^{-k}\)
in deterministic bit time polynomial in \(n,b,k\),
including all query and answer bits.
For a nonexpansive map (\(\gamma=0\) allowed
in the Lipschitz condition), we can instead return a dyadic \(x\in C\)
with residual
\(\|x-f(x)\|_\infty\le 2^{-k}\),
in time polynomial in \(n,k\).
\end{theorem}

For a contraction, a residual bound \(e\) also gives the distance
bound \(e/\gamma\). Shellman and Sikorski~\cite{Sik} obtained
polylogarithmic dependence on inverse accuracy in fixed dimension,
with an exponent growing with the dimension. Chen, Li and
Yannakakis~\cite{CLY} balanced the remaining possible locations to
obtain a residual of at most \(\epsilon\) in
\(O(n\log(1/\epsilon))\) exact evaluations. Computing the query
points incurs an additional cost. Feodorov and Haslebacher~\cite{FH}
combined faster balancing with dimension splits~\cite{CLYsp}
to obtain subexponential arithmetic time in \(n+\log(1/\epsilon)\)
with exact evaluations. See also the subsequent work on
decomposition~\cite{CLYdec}.\footnote{For other norms see~\cite{STWon,HKS,HLSW,DFZ},
for monotone maps see~\cite{BFGMS}, and for the complexity of total
search see~\cite{DP,FGMS}.} Soleymani and Farina~\cite{SF} obtained
quasipolynomial arithmetic complexity for approximating the location,
counting exact evaluations and using logarithms of inverse gap and
inverse error as parameters. Theorem~\ref{thm:prices} combines
balancing with recursive amplification of extremal labels.

\begin{proof}
First assume \(\gamma=2^{-b}>0\) and set
\[
 \epsilon=2^{-k},\qquad \delta=\gamma\epsilon/16,\qquad R=2/\delta .
\]
Extend \(f\) to \(\RR^n\) by clipping each input coordinate to
\([-1,1]\). This preserves the Lipschitz bound and the fixed point.
Run Theorem~\ref{thm:prices} with \(N=n\) and demand bound \(R\).
At a queried price vector \(p\), set \(x=\delta p\).
To obtain a label, clip \(x\) to \(C\), round it to the nearest
\(\delta\)-grid point, and request \(b+k+4\) bits from the oracle.
Subtract the original \(x\) from the answer and return the signs,
breaking ties arbitrarily. Input rounding contributes at most
\(\delta/2\) to the evaluation error and the oracle at most
\(\delta\). Thus the labels recover the sign of every component of
\(f(x)-x\) with magnitude greater than \(2\delta\).

For the analysis, the fixed point defines one hypothesis with both
slots \((i,h)\), \(h=\pm1\), at every coordinate. Assign the demand
\[
 a_{i,h}=\left\lfloor\frac{h v_i-\epsilon/2}{\delta}\right\rfloor.
\]
These integers lie in \([-R,R]\). The algorithm need not know them.
Writing \(d_{i,h}(p)=a_{i,h}-hp_i\), we have
\[
 h(v_i-x_i)-\delta
 <\delta d_{i,h}(p)+\epsilon/2
 \le h(v_i-x_i).
\]
Suppose that \((i,h)\) maximizes the surplus and its surplus is
nonnegative. With \(u=h(v_i-x_i)\) the displayed bounds give
\(u\ge\epsilon/2\) and \(\|v-x\|_\infty\le u+\delta\).
Since \(f(v)=v\), the movement of coordinate \(i\) satisfies
\[
 h(f_i(x)-x_i)\ge u-(1-\gamma)(u+\delta)
     \ge \gamma\epsilon/2-\delta=7\delta .
\]
Even after rounding, the reported sign is \(h\).
This proves consistency of the hypothesis.

At the returned prices all its surpluses are nonpositive, so the
corresponding \(x\) satisfies
\(\|x-v\|_\infty\le\epsilon/2+\delta\). Clipping and rounding to
the same grid yield a dyadic point within \(\epsilon\) of \(v\).
The running time is polynomial in \(n,b,k\) because
\(\log R=b+k+5\), and every oracle input and answer is a string
of polynomial length.

For a nonexpansive \(f\), apply the contraction result to
\(g=(1-\epsilon/2)f\) with target distance \(\epsilon/4\)
as in~\cite{CLY,SF}. At the fixed point \(v_g\), the residual
\(\|v_g-f(v_g)\|_\infty\) is at most \(\epsilon/2\) because
\(f(v_g)\in C\). Moving to the returned point changes this residual
by at most \(2(\epsilon/4)\), by nonexpansiveness. To simulate an
\(m\)-bit query to \(g\), request \(m+1\) bits from \(f\),
multiply by \(1-\epsilon/2\), and round to the nearest \(m\)-bit
dyadic. The total error is at most \(2^{-m}\).
\end{proof}

\subsection{Exact piecewise-affine computation}

For an explicitly given piecewise-affine map, a small residual
identifies a region on which an exact fixed point can be found by
LP. Consider a circuit with \(n\) inputs and \(n\) outputs, built
from scalar affine, minimum and maximum gates~\cite{EY,FGMS}.
A gate output may feed several later gates. A \emph{selection}
specifies one attaining operand at each minimum or maximum.
Adding a variable for each gate output gives linear constraints for
the inputs and outputs realizing a selection. Ties are allowed, so
each such cell is closed.

\begin{corollary}[Exact fixed points for rational circuits]\label{cor:pa}
Given such a circuit with binary rational coefficients, a
deterministic algorithm uses polynomial bit time to return either an
exact fixed point in \(C\) or a counterexample to the nonexpansive
self-map property. A counterexample consists of one point mapped
outside \(C\) or two points violating nonexpansiveness. All points
returned are rational.
\end{corollary}

\begin{proof}
Fix a selection and minimize \(e\) over the realizing cell in \(C\),
subject to \(-e\le f_i(x)-x_i\le e\) for every \(i\). Every such LP
has encoding length bounded by a common polynomial in the circuit
length. By the usual vertex bounds~\cite{GLS}, its minimum, if
positive, is at least the reciprocal of an integer of polynomial
bit length. Choose \(k\), polynomial in the input length, so that
\(2^{-k}\) is strictly smaller than this lower bound. Every cell
realizing a residual at most \(2^{-k}\) then contains an exact fixed
point. This precision argument is due to Etessami and
Yannakakis~\cite{EY}.

Run Theorem~\ref{thm:contraction}
for this \(k\), in its nonexpansive version.
Evaluate the circuit exactly before rounding each oracle reply.
With the input and a realizing selection prescribed, the gate values
are the unique solution of a rational linear system. Vertex bounds
therefore give polynomial bit length for all reduced gate values
on short-bit inputs.
Include evaluation at the point output by that theorem.
If an evaluation breaks range or violates nonexpansiveness
upon comparison with past evaluations, stop
and return this counterexample.

If all pass, these exact evaluations extend to
a nonexpansive self-map of the cube by McShane extension
\cite{McSh}
(coordinatewise, take the minimum of recorded
values plus sup norm distance to the corresponding sample,
then clip).
So the residual guarantee holds, even for the original circuit, at
the output point.
This extension argument for total search
(cf.~\cite{CLY})
applies also to each prefix of passing evaluations,
so the time bound holds in any case.
Thus we need only make a selection at the output point
and solve linear feasibility for a fixed point in that cell
by bit-polynomial LP.
\end{proof}

\subsection{Offsets and asymptotics}

A discounted game gives affine gates whose slopes contain
transition probabilities and discount factors, while rewards enter
as offsets. The offsets may be arbitrary reals. To study vanishing
discount, we also need to recover a single selection valid throughout
an interval \(0<t<t_*\). We do both by determining which linear
constraints are tight. Their determinant tests depend linearly on
the offsets, permitting Frank--Tardos rounding in the spirit of
Tardos's algorithm for LP~\cite{Tardos}.

\begin{proposition}[Real offsets and a vanishing parameter]\label{prop:asym}
Allow affine gates in Corollary~\ref{cor:pa}
with operands \(y_j\) and output
\[
 z=\sum_j a_j(t)y_j+b(t),
\]
where the polynomials are given by coefficient lists. The slope
polynomials \(a_j\) have binary rational coefficients, and the offset
polynomial \(b\) has arbitrary real coefficients.
There is a deterministic strongly polynomial algorithm returning
either a fixed point in \(C\) or a counterexample to the nonexpansive
self-map property as in Corollary~\ref{cor:pa}. The coordinates
are rational functions of \(t\), represented by ratios of real
polynomials, valid for all sufficiently small \(t>0\).
With a fixed point we can also return a selection, valid in the
same sense.
Here degrees and slope encoding lengths count as part
of discrete input size.
\end{proposition}

This includes nonparametric input (and then the output
can be taken constant).

\begin{proof}
We solve on short-bit data and transfer the result back.
For this we recognize possible outcomes by determinant signs.

Indeed, at fixed \(t\), use evaluation constraints with
selections prescribed. Then asking for a fixed point in \(C\)
gives linear conditions. To ask if the outputs \(z,z'\)
at points \(x,x'\in C\) can violate nonexpansiveness,
specify an output coordinate \(i\) and sign \(h\), and consider
\[
  h(z_i-z'_i)\ge \tau+u,\qquad
  \|x-x'\|_\infty\le\tau\le2,\qquad 0\le u\le1 .
\]
A solution with \(u>0\) witnesses a violation. For a range violation
prescribe a single selection, coordinate \(i\) and sign \(h\), and
impose \(hz_i\ge1+u\) with \(0\le u\le1\). Again success means
\(u>0\).

With only its weak inequalities, each system describes a polytope,
since the inputs lie in the cube and the gate outputs are bounded
there. A successful outcome thus has a vertex witness, obtained by
maximizing \(u\) when a violation is sought. A square invertible
system of tight rows determines this vertex. Its determinant and
the bordered determinants formed with each inequality row and the
right-hand side column determine whether the vertex satisfies
the inequalities. They also determine the sign of \(u\).

All these determinants are univariate polynomials with
coefficients linear in the vector \(w\)
of all offset coefficients together with a constant~\(1\)
(use the \(1\) entry also for tests independent of offsets).
After a common positive denominator is cleared, these linear
forms have integer coefficients of \(\ell_1\)-norm bounded
by \(2^{\poly(L)}\), uniformly over selections,
with \(L\) our input size counting each offset coefficient as one entry.
Indeed, clear all slope denominators by their product,
then use the determinant expansion
(degrees polynomial, at most polynomially many factors per term,
at most \(2^{\poly(L)}\) terms).
Offsets appear only in the right-hand side column.

By Lemma~\ref{lem:ft} use short-bit offsets (normalizing
the rounded vector so the constant is again one)
preserving all these coefficient signs.
At the new offsets, use a positive
\(t_0=2^{-\poly(L)}\)
small enough that the lowest nonzero coefficient
decides every determinant sign at \(t_0\)
whenever the determinant is nonzero as a polynomial.
After clearing the new denominators, its coefficients are integers
of polynomial bit length. A geometric series bounds the terms
after the lowest nonzero one, so this sign agrees with that at all
sufficiently small positive \(t\). Identically zero determinants
stay zero.

Solve the resulting short-bit problem (new offsets, \(t=t_0\))
by Corollary~\ref{cor:pa}.
Recover a successful vertex and defining tight rows
as above, by ordinary bit-polynomial LP
(use selections at the obtained points, maximizing \(u\)
where applicable, then getting a vertex).
The same determinant signs hold with original offsets
at all sufficiently small positive \(t\).
Return the corresponding vertex by Cramer's rule,
with the numerator and denominator polynomials formed by Lagrange
interpolation. The degree bounds above allow interpolation from
evaluations at polynomially many small integers, avoiding polynomial
division.
This is strongly polynomial by Section~\ref{sec:arithmetic}
(interpolation uses linear combinations with short-bit coefficients).
\end{proof}
